\documentclass[a4paper,10pt]{article}
\usepackage[left=1.6cm,right=1.6cm,top=1.4cm,bottom=1.4cm]{geometry}
\usepackage[fontset=fandol]{ctex}
\usepackage{enumitem}
\usepackage{titlesec}
\usepackage{xcolor}
\usepackage{hyperref}
\usepackage{tabularx}
\pagestyle{empty}
\definecolor{accent}{HTML}{111827}
\hypersetup{colorlinks=true,urlcolor=accent,pdfauthor={Resume Owner},pdftitle={Chinese Resume Template}}
\titleformat{\section}{\normalsize\bfseries\color{accent}}{}{0pt}{}[\titlerule]
\titlespacing*{\section}{0pt}{5pt}{2pt}
\setlist[itemize]{topsep=1pt,itemsep=1pt,parsep=0pt,partopsep=0pt}

\newcommand{\cnheading}[4]{
  \item\begin{tabularx}{\linewidth}{@{}Xr@{}}
  \textbf{#1} & #4 \\
  \textit{#2} & \textit{#3}
  \end{tabularx}\vspace{-3pt}
}
\newcommand{\cnitem}[1]{\item\small #1}

\begin{document}
\begin{center}
  {\LARGE\bfseries 你的姓名}\par\vspace{4pt}
  \small
  \href{mailto:you@example.com}{you@example.com}\quad\textbar\quad
  +1 (555) 010-1234\quad\textbar\quad
  \href{https://github.com/your-handle}{github.com/your-handle}
\end{center}

\section{教育背景}
\begin{itemize}[label={},leftmargin=0pt]
\cnheading{示例大学}{计算机科学硕士}{2024.08--2026.05}{示例城市}
\cnheading{样例学院}{计算机工程学士}{2020.08--2024.05}{样例城市}
\end{itemize}

\section{实习经历}
\begin{itemize}[label={},leftmargin=0pt]
\cnheading{示例实验室}{软件工程实习生}{2025.05--2025.08}{示例城市}
\begin{itemize}[leftmargin=1.3em]
\cnitem{构建日均处理 \textbf{120 万条事件}的生产服务，将 p95 延迟降低 \textbf{38\%}。}
\cnitem{增加集成测试与部署检查，将每季度回滚事故从 5 次减少到 1 次。}
\end{itemize}
\end{itemize}

\section{项目经历}
\begin{itemize}[label={},leftmargin=0pt]
\cnheading{QueueScope}{Python / FastAPI / PostgreSQL / Docker}{2025.01--2025.04}{}
\begin{itemize}[leftmargin=1.3em]
\cnitem{为异步任务构建运维面板，支持链路检索、失败重试与基于角色的权限控制。}
\cnitem{针对 \textbf{5,000 个并发任务}增加负载测试，并记录系统性能边界。}
\end{itemize}
\cnheading{FieldNotes}{TypeScript / React / SQLite}{2024.09--2024.12}{}
\begin{itemize}[leftmargin=1.3em]
\cnitem{为结构化笔记实现离线优先同步与确定性冲突解决。}
\end{itemize}
\end{itemize}

\section{专业技能}
\begin{itemize}[leftmargin=1.3em]
\cnitem{\textbf{编程语言：}Python、TypeScript、SQL}
\cnitem{\textbf{框架与基础设施：}FastAPI、React、PostgreSQL、Docker、GitHub Actions}
\end{itemize}
\end{document}
